\section{The Associative Engine}

An earlier section showed that the bulk holds a database, balanced by its geometry, and that the balance is paid for by the curvature itself. This section asks a sharper question. Not whether the mesh can store, but whether it can \emph{remember} the way a mind remembers. A mind does not recall by address. It does not say "fetch the word at slot $40{,}912$." It recalls by content. A fragment of a song, a half-seen face, the first three letters of a name, and the whole returns. That is the mark of an associative memory, and this section measures the mesh against it.

\begin{definition}[Content-addressable memory]
A \emph{content-addressable memory} recovers a whole stored word from a corrupted or partial piece of it. The query is not a location but a fragment of the content, and the memory returns every stored word that matches.
\end{definition}

The claim of this section is that the $\{3,4,3,4\}$ bulk is such a memory, and that it is an exceptionally good one, because hyperbolic space lays out an associative memory exponentially better than flat space does. Every number that follows was \emph{measured}, each against a control that could have failed.

\subsection{The headline}

One sentence carries the result. A $24$-site $\{3,4,3,4\}$ bulk holds $3000$ stored words within search radius $3$, and that radius does \emph{not grow} when the memory quadruples.

Three facts ride alongside it.

\begin{itemize}
\item The content search is exact. It recovers stored words from corrupted cues.
\item The same engine runs unchanged on all $42$ buildable tessellations of the catalog.
\item The whole machine runs on the GPU. In one dispatch of $1.14$ milliseconds it searches $200{,}000$ docks at once.
\end{itemize}

\begin{result}[Associative memory is hyperbolic]
Content-addressable memory is exponentially better laid out on hyperbolic space than on flat space. The search radius that covers the whole memory grows logarithmically with the number of stored words on the bulk, and polynomially on a flat cubic memory of the same size.
\end{result}

\subsection{Why the bulk is the right place}

The reason is geometric and it is the whole point of the section. A memory is useful only if a query can reach all of it quickly. So the question is how far a query has to travel to touch every stored word.

A flat cubic memory of $N$ docks has diameter about $N^{1/3}$. To broadcast a query to every stored word, the activation wave must cross that diameter. The latency therefore grows \emph{polynomially} in the size of the memory. Double the memory often enough and the search slows without bound.

The bulk has \emph{logarithmic} diameter. This is the geometry result established earlier, that a ball of radius $r$ in $\{3,4,3,4\}$ holds exponentially many docks, so the radius needed to enclose $N$ docks is only of order $\log N$. The same broadcast that crawls across a flat memory covers the whole bulk in $O(\log N)$ beats.

\begin{principle}[The logarithmic-diameter win]
The hyperbolic diameter is the win. A query reaches everything in logarithmic time because the bulk has no far-away docks. Every dock sits within $O(\log N)$ beats of every other, so there is no remote corner for a stored word to hide in.
\end{principle}

This is measured on both the CPU and the GPU, against a flat cubic control, in the deepest experiments of the suite. The contrast is not asserted from the geometry alone. It is checked dock for dock.

\subsection{The depth ledger}

The experiments are graded on a depth rubric. An L1 experiment shows a thing works once. An L2 experiment shows it degrades or scales sensibly. An L3 experiment compares the bulk against a flat control and reports the \emph{scaling}, not a single point, so that a wrong theory would produce the wrong slope. The load-bearing claims of this section are all L3.

\begin{table}[ht]
\centering
\begin{tabular}{@{}lll@{}}
\toprule
experiment & depth & measured result \\
\midrule
search latency & L3 & bulk radius $3$ at $750$ and $3000$ docks ($\Delta 0$), flat $\Delta 9$ \\
parallel cost & L3 & cost constant ($31$) over an $8\times$ size rise, bulk $\Delta 1$, flat $\Delta 18$ \\
spreading activation & L3 & wave covers $1500$ docks in $3$ beats, arrival equals graph distance \\
capacity vs curvature & L3 & growth ratio rises with curvature $1.66 \to 2.16$, radius falls $11 \to 8$ \\
Hopfield recall & L3 & emergent layer recalls $1.0$, bare knit recalls $0.0$ (chance $0.167$) \\
\bottomrule
\end{tabular}
\caption{The five deep, control-backed measurements of the associative engine.}
\label{tab:associative-deep}
\end{table}

Read the first row carefully, because it is the heart of the section. On the bulk the coverage radius stays at $3$ as the memory grows from $750$ to $3000$ docks. The increment is zero. It reaches $5$ only at $200{,}000$ docks. A flat cubic memory's radius grows by $9$, then by $18$, for the same increases. The radius increment under a fixed size jump is the logarithm in disguise, and the bulk wins it decisively.

\begin{result}[Logarithmic search latency]
On the bulk, the radius that covers the whole memory grows logarithmically with the number of stored words. From $750$ to $3000$ docks the radius does not grow at all. On a flat cubic memory the same radius grows polynomially, by $9$ then $18$. The parallel-match cost is constant in size on the bulk, while the communication radius grows logarithmically there and polynomially on the flat memory.
\end{result}

\subsection{The parallel match}

The simplest query mode is a single shout. Broadcast the cue to \emph{every dock at once}, and let each dock decide locally whether it matches. There is no scan, no index walk, no central coordinator. Every dock holds a stored word, the cue arrives, each dock compares its own word against the cue in parallel, and the matching docks light up.

On the GPU this is exactly one dispatch.

\begin{table}[ht]
\centering
\begin{tabular}{@{}p{0.42\textwidth}p{0.5\textwidth}@{}}
\toprule
what it does & measured result \\
\midrule
the parallel match, one dispatch over all docks & the GPU responder set equals the CPU ground truth, exact and partial. $200{,}000$ docks searched in $1.14$ ms in one dispatch \\
\bottomrule
\end{tabular}
\caption{The parallel match, verified bit for bit against the CPU.}
\label{tab:parallel-match}
\end{table}

The verification matters more than the speed. The GPU result is not trusted on its own. The set of docks the GPU reports as matching is compared against the set the CPU computes by exhaustive search, and the two sets are equal, for both exact matches and partial matches. The speed is then real, $200{,}000$ docks in just over a millisecond, because the answer is also correct.

\subsection{The broadcast wave}

The second query mode is a wave. The cue activates a starting region, and activation spreads outward shell by shell, one ring of docks per beat. This is exactly what cognitive science calls \emph{spreading activation}, the process by which thinking of one idea makes nearby ideas easier to reach.

\begin{table}[ht]
\centering
\begin{tabular}{@{}p{0.42\textwidth}p{0.5\textwidth}@{}}
\toprule
what it does & measured result \\
\midrule
the query wave, one ring per beat & the GPU arrival beats equal the CPU breadth-first search dock for dock, zero mismatches. $200{,}000$ docks covered in $5$ beats, $3.15$ ms \\
\bottomrule
\end{tabular}
\caption{The spreading-activation wave, matched against a breadth-first search.}
\label{tab:wave}
\end{table}

Two things about the wave are worth drawing out. First, because the bulk diameter is logarithmic, the wave reaches the entire memory in a handful of beats. Five beats for $200{,}000$ docks. The same wave on a flat cubic memory would have to cross a polynomial diameter and would take far longer. Second, the arrival time of the wave at any dock \emph{equals its graph distance} from the cue. The wave is a beat-by-beat breadth-first search, and the GPU reproduces the CPU breadth-first search exactly.

\begin{principle}[Spreading activation is the query wave]
Spreading activation through a semantic memory is the bulk query wave, made logarithmic by hyperbolic geometry. The wave's arrival time at a dock is its distance from the cue, and the geometry compresses every distance to $O(\log N)$.
\end{principle}

\subsection{Exponential capacity per radius}

The third deep fact answers a storage question. How much can be stored in a ball of a given radius? The bulk's answer is what makes it the right substrate.

\begin{table}[ht]
\centering
\begin{tabular}{@{}lll@{}}
\toprule
experiment & depth & measured result \\
\midrule
capacity scaling & L2 & bulk per-radius growth ratio $13.6$ versus flat cubic $1.10$ \\
capacity vs curvature & L3 & growth ratio rises with curvature $1.66 \to 2.16$, latency falls $11 \to 8$ \\
\bottomrule
\end{tabular}
\caption{Capacity per radius, exponential on the bulk and tied to curvature.}
\label{tab:capacity}
\end{table}

The per-radius growth ratio is $13.6$ on $\{3,4,3,4\}$ against $1.10$ flat. A ball of fixed radius in the bulk holds exponentially more stored words than the same-radius ball in flat space. That is the exponential room of hyperbolic geometry, now read as memory capacity.

Capacity also tracks curvature. Across the catalog of tessellations the growth ratio rises with curvature, from $1.66$ to $2.16$, while the coverage radius falls, from $11$ to $8$. More curvature buys both. More words fit in a given ball, and the search to reach them is shorter.

\begin{result}[Exponential capacity, tied to curvature]
A ball of fixed radius in the bulk holds exponentially more stored words than the same ball in flat space, with a per-radius growth ratio of $13.6$ against $1.10$. Across tessellations capacity rises and latency falls as curvature increases. Curvature buys capacity and speed together.
\end{result}

\subsection{Robust content-addressable recall}

The deep scaling results rest on a recall core that works and degrades gracefully. These are the experiments that confirm the memory does what a memory must.

\begin{table}[ht]
\centering
\begin{tabular}{@{}lll@{}}
\toprule
experiment & depth & measured result \\
\midrule
exact recall & L1 & recall $1.0$, false positives $0$, on $1500$ docks \\
numeric search & L1 & max, min, next, previous all correct, cost $17$ passes, constant in dock count \\
noisy recall & L2 & recall $1.0$ clean, $1.0$ at $10\%$, $1.0$ at $20\%$, $0.6$ at $40\%$ corruption \\
graded recall & L2 & full-cue fidelity $1.0$, smooth decay as the cue shrinks (max step $0.28$) \\
interference & L2 & well-separated words recall $1.0$, overlapping words crosstalk to $0$ \\
\bottomrule
\end{tabular}
\caption{The recall core, exact and graceful.}
\label{tab:recall}
\end{table}

Exact recall is perfect, with no false positives, on $1500$ docks. Noisy recall holds at full fidelity up to twenty percent corruption of the cue, then declines, reaching $0.6$ at forty percent. Graded recall degrades smoothly as the cue is whittled down, with no fidelity step larger than $0.28$, provided the stored word is distributed across the dock's sites rather than packed into a few. Interference behaves as a real memory should. Well-separated words recall cleanly, and only deliberately overlapping words crosstalk.

\subsection{Human-style memory on the base layer}

The engine reproduces the textbook properties of human memory, and it does so on the base layer of the model, with no special machinery added.

\begin{table}[ht]
\centering
\begin{tabular}{@{}lll@{}}
\toprule
experiment & depth & measured result \\
\midrule
priming & L2 & unprimed query $3$ beats, primed query $1$ beat, $2$ beats saved \\
chaining & L2 & a free-association chain visits $12$ distinct related docks without collapsing \\
sequence recall & L2 & an $8$-step stored sequence replays fully in order from the first cue \\
tessellation sweep & L2 & exact recall $1.0$ on all $42$ buildable tessellations \\
dense vs classical capacity & L2 & classical capacity $14$ ($0.12N$), dense capacity $120$ ($1.0N$) \\
\bottomrule
\end{tabular}
\caption{Human-style memory behaviors, realized on the base layer.}
\label{tab:human-memory}
\end{table}

Priming is real. A query that follows a related cue completes in one beat instead of three, two beats saved, because the earlier cue left the neighborhood warm. Free association chains. Starting from one cue, a walk visits twelve distinct content-related docks without collapsing back on itself. A stored sequence of eight steps replays fully in order from its first cue alone.

The sweep is the universality claim. All $42$ buildable tessellations recall perfectly with one code path. Recall depends on the dock graph itself, not on which tessellation built it. The geometry sets only the capacity and the latency.

\begin{principle}[Universal across the catalog]
The associative engine is universal across the tessellation catalog. All $42$ buildable tessellations recall with one unchanged code path. Recall is a property of the dock graph. The geometry tunes capacity and latency, but it does not gate whether the memory works.
\end{principle}

\subsection{The dense Hopfield layer}

A Hopfield network is an attractor memory. Corrupt a stored pattern and let the network relax, and it slides downhill in energy until it settles on the nearest clean pattern. The dense, modern variant of this layer is an emergent structure on the mesh, and it cleans up cues that the bare search cannot.

\begin{table}[ht]
\centering
\begin{tabular}{@{}p{0.42\textwidth}p{0.5\textwidth}@{}}
\toprule
what it does & measured result \\
\midrule
the dense (emergent) Hopfield recall & a $0.648$-corrupted cue recalled to its prototype at $1.000$, GPU equals CPU $100\%$. Benchmark $4096$ neurons, $200$ patterns, $5$ steps, $7.37$ ms, recall $1.000$ \\
\bottomrule
\end{tabular}
\caption{The dense Hopfield attractor layer, an emergent clean-up stage.}
\label{tab:hopfield}
\end{table}

A cue corrupted nearly two-thirds of the way to noise is pulled back to its prototype with full fidelity. The relaxation of $4096$ neurons over $200$ patterns runs in under eight milliseconds on the GPU, bit for bit equal to the CPU.

\subsection{The honest findings}

Four findings were discovered rather than designed, and reporting them is part of the method. A measured result that comes out the wrong way teaches something, and these did.

\begin{result}[Graceful recall is representation, not geometry]
The smooth degradation of recall is a property of the word's representation, not of the curvature. A dense default word packs identity into the low sites and gives a sharp information-capacity cliff rather than a smooth decay. A \emph{distributed} word, with identity spread across all sites of the dock, is what yields human-like graded recall. The experiment says so explicitly, and the cliff returns if the representation is changed back.
\end{result}

\begin{result}[Attractor recall is emergent, not base]
The bare reversible knit cannot do attractor recall. It conserves charge and has no energy to descend, so it cleans up at chance, recall $0.0$ against a chance of $0.167$, where the emergent dissipative Hopfield layer recalls at $1.0$. Content-addressable search is a base-layer ability. Attractor clean-up is an emergent-layer ability. This split is demonstrated, not assumed.
\end{result}

The third finding is a caution against overclaiming. At small sizes the bulk's absolute radius can exceed the flat one, because the $\{3,4,3,4\}$ growth ratio is a modest $1.575$ asymptotically. The logarithmic advantage lives in the \emph{scaling}, the radius increment under a size increase, not in the absolute radius at a few thousand docks. The experiments assert the scaling, which is the honest claim.

The fourth finding bounds the sequence demonstration. A greedy walk that visits only unvisited neighbours dead-ends after eight docks at a memory of $1500$, because the bulk's short simple paths run out. So the demonstrated sequence is length eight, replayed fully in order, and the bound is flagged rather than hidden.

\subsection{Full-suite confirmation}

All sixteen associative experiments pass. Run as part of the whole repository gate after the build, the suite returns no crashes and no new regressions. The known pre-existing honest-negatives elsewhere in the model are unchanged. Nothing in the associative engine broke anything, and nothing was tuned to make a number come out right.

\begin{principle}[The GPU is the practical machine]
The match is one dispatch, $200{,}000$ docks in $1.14$ ms. The wave covers the same in $5$ beats. The dense Hopfield relaxes $4096$ neurons in $7.37$ ms. Each result is bit for bit equal to the CPU ground truth, so the speed is honest and the engine is a practical machine, not a thought experiment.
\end{principle}

The bulk does not merely store. It remembers by content, the way a mind does, and it remembers exponentially better than flat space because its diameter is logarithmic. Spreading activation is the query wave. Priming saves beats. Free association chains. Recall is graded. Sequences replay in order. The memory is hyperbolic by construction, and the measurements say the construction works.
